\section{Related Work}
\label{sec:related_work}

The field of industrial fault diagnosis has evolved from signal processing methods to advanced deep learning architectures. However, the reliance of these models on large-scale, labeled training data remains a significant barrier for practical deployment. In this section, we review the existing literature in four areas related to our work: (1) Zero-Shot Fault Diagnosis (ZSFD) in Industry; (2) Generative Adversarial Networks (GANs) and Variational Autoencoders (VAEs) for industrial signals; (3) recent developments in Denoising Diffusion Probabilistic Models (DDPM); and (4) the emerging field of cross-modal semantic alignment via large-scale pre-trained models.

\subsection{Zero-Shot Fault Diagnosis (ZSFD) in Industry}
Zero-Shot Learning (ZSL) has emerged as an innovative solution to the data scarcity problem in industrial monitoring, focusing on classifying unseen categories present in the test set but absent during training. Early ZSFD methods relied on binary attributes—hand-crafted by domain experts—to bridge the gap between seen and unseen mechanical faults. For example, a bearing fault might be described by a set of attributes such as "high-frequency," "periodic," and "impact-oriented." Traditional ZSL frameworks like ALE, SJE, and DEVISE formulated the problem as a linear or non-linear mapping between raw features and semantic attributes. However, in industrial settings, these models often suffer from the "hubness problem," where a few semantic prototypes become near neighbors to many unrelated signal features in high-dimensional space.

Recent applications of ZSL in industrial domains have expanded into complex systems. For instance, in high-voltage circuit breakers (HVCBs), cross-domain ZSL has been used to transfer knowledge from laboratory environments to real-world power plants, where fault data for actual field failures is virtually non-existent. Similarly, in wind turbine gearbox monitoring, researchers have utilized label embeddings to represent various mechanical states without requiring explicit samples of every rare failure mode. These specialized applications often employ three phases: feature extraction from raw signals (using 1D CNNs or LSTM), label embedding into a semantic space (using Word2Vec or domain knowledge), and statistical feature embedding to align the two domains. Despite these advances, traditional ZSL often struggles with the "domain shift" problem, where the statistical distribution of synthesized features does not perfectly match real-world sensor drift, especially in complex chemical processes like the Tennessee Eastman Process (TEP). Furthermore, the reliance on static word vectors like Word2Vec fails to capture the technical nuances of industrial manuals. Our work addresses these issues by leveraging a pre-trained semantic manifold that provides a more continuous and objective representation of the fault space.

\subsection{Generative Models: GANs and VAEs}
To improve the robustness of ZSFD, generative models like Generative Adversarial Networks (GANs) and Variational Autoencoders (VAEs) have been developed to "hallucinate" synthetic features for unseen fault classes. For instance, f-CLSWGAN was one of the first models to use a Wasserstein-GAN to generate vibration features conditioned on semantic attributes. Other variants, such as the Generative Dual Adversarial Network (GDAN), utilize a discriminator to estimate the similarity level with respect to each visual-textual (or signal-textual) feature pair, using a dual adversarial loss function to ensure class-specific generation. In cases of imbalanced data, GANs have been shown to capture distribution information by fusing information from multiple domains.

While GANs are computationally efficient, they are notoriously difficult to train, frequently suffering from mode collapse or unstable gradients. VAE-based approaches, such as CADA-VAE (Cross-Alignment Disentangled Auto-Encoder), offer a more stable probabilistic framework by learning shared cross-modal latent features. However, VAEs tend to produce "blurry" synthetic features that lack the fine-grained temporal details of real vibration signals. In industrial contexts, where high-frequency transients and subtle periodic impulses are critical for identifying early-stage bearing wear, the limitations of GANs and VAEs often lead to poor generalization. Specifically, they often fail to capture the intricate cross-sensor correlations necessary for accurate system-wide diagnosis. Hybrid VAE-GAN models have been explored but still struggle to model the complex, non-Gaussian sensor distributions found in modern chemical processes.

\subsection{Diffusion Models for Industrial Synthesis}
Denoising Diffusion Probabilistic Models (DDPM) have recently emerged as the state-of-the-art for high-fidelity generative tasks. By iteratively refining a noise vector into a clean sample through a multi-step denoising trajectory, diffusion models provide a more stable and expressive generative framework than GANs. In the industrial realm, diffusion models are beginning to show superiority in synthesizing realistic vibration data and chemical process time-series. Unlike GANs, diffusion models optimize a variational lower bound that allows for better coverage of the entire data manifold. This prevents the "missing mode" problem and ensures that even rare fault characteristics are represented in the generated samples. Furthermore, the iterative nature of diffusion allows for more flexible conditioning, making it an ideal candidate for combining diverse semantic inputs with physical signal distributions. Recent advances in guided and latent diffusion (LDM) have reduced the computational overhead by performing the denoising process in a compressed latent space, which is particularly relevant for high-sampling-rate industrial signals.

\subsection{Cross-Modal Alignment and CLIP}
Models like CLIP (Contrastive Language-Image Pretraining) have demonstrated unprecedented success in aligning visual concepts with natural language by training on billions of image-text pairs. This allows CLIP to learn a joint embedding space where semantically similar concepts are geometrically close. In industrial fault diagnosis, this offers a breakthrough opportunity to replace limited, hand-crafted attributes with rich descriptions. For example, a description like "inner race pitting at 1000 RPM with low lubrication" carries significantly more structured information than a simple class index.

Industrial CLIP variants have been proposed to align vibration spectrograms with failure mode analysis (FMA) reports. Our CESM-Diff framework is the first to integrate a dual-path diffusion process with the CLIP semantic manifold for industrial sensing. This allows for the synchronization of the feature manifold with a context-rich, pre-trained linguistic space, enabling the framework to generalize to novel fault categories with higher fidelity and semantic grounding. Furthermore, by framing the diagnosis task as a text-to-signal matching problem, we can exploit the inherent zero-shot capabilities of the CLIP manifold, allowing the system to handle new categories described only in plain text, thus significantly reducing the manual overhead of domain-specific feature engineering.
